% ══════════════════════════════════════════════════════════════════════════
% §2  Hardware Reality: Why Your O(n³) is Slow
% ══════════════════════════════════════════════════════════════════════════
\section{Hardware Reality: Why Your \texorpdfstring{$O(n^3)$}{O(n³)} is Slow}
\label{sec:hardware}

In C++, \textbf{memory is the bottleneck, not arithmetic}.
A modern CPU can execute $\sim\!10^{10}$ floating-point operations per second,
but fetching a cache line from main RAM takes $\sim\!100$ ns---long enough
for $\sim\!200$ multiplies to have completed.

\subsection{Cache Locality}

CPUs load data in \emph{cache lines} (typically 64 bytes = 8 doubles).
If your next access is a neighbour in memory, it is already in cache
(\textbf{cache hit}).  If you jump across rows or columns, the CPU stalls
while fetching a new line (\textbf{cache miss}).

\begin{definition}[Row-major vs.\ Column-major]
In \textbf{row-major} layout (C/C++ default), the element
$A_{i,j}$ is stored at offset $i \cdot n + j$.
Incrementing $j$ visits contiguous memory.

In \textbf{column-major} layout (Fortran, MATLAB, \textbf{Eigen default}),
the element $A_{i,j}$ is stored at offset $j \cdot m + i$.
Incrementing $i$ visits contiguous memory.
\end{definition}

\begin{warningbox}
Eigen stores matrices in \textbf{column-major} order by default.
If you write a row-traversal loop \texttt{for(i) for(j) A(i,j)},
every inner-loop access jumps by $m$ doubles---a cache miss every time.
Swap the loop order, or use \texttt{Eigen::RowMajor}.
\end{warningbox}

The following benchmark makes the difference concrete:

\lstinputlisting[caption={Cache locality: row-major vs.\ column-major traversal.},
                 label={lst:cache}]
                {code/cache_locality.cpp}

\subsection{SIMD: Single Instruction, Multiple Data}

Modern x86 CPUs have \textbf{AVX-2} (4 doubles/cycle) or \textbf{AVX-512}
(8 doubles/cycle) vector units.  When your data is \emph{contiguous and
aligned}, the compiler auto-vectorises your loop and processes 4--8 elements
per instruction.  When data is scattered, the vector unit idles.

\begin{tipbox}
You almost never write SIMD intrinsics by hand.
The rule is simpler: \textbf{keep your data contiguous}, and the
compiler (with \texttt{-O2 -march=native}) will vectorise for you.
Libraries like Eigen and MKL are already written to exploit this.
\end{tipbox}

\subsection{The Takeaway}

\begin{center}
\renewcommand{\arraystretch}{1.3}
\begin{tabular}{@{}lcc@{}}
\toprule
\textbf{Access pattern} & \textbf{Cache behaviour} & \textbf{Relative speed} \\
\midrule
Contiguous (correct order) & Hits    & $1\times$ \\
Stride-$n$ (wrong order)   & Misses  & $3$--$10\times$ slower \\
Random access              & Misses  & $10$--$50\times$ slower \\
\bottomrule
\end{tabular}
\end{center}
